\documentclass[11pt]{article}
\usepackage[a4paper,margin=2.4cm]{geometry}
\usepackage{fontspec}
\usepackage{unicode-math}
\usepackage{xcolor}
\usepackage{booktabs}
\usepackage{hyperref}
\setmainfont{SourceSerif4}[Path=fonts/,Extension=.otf,UprightFont=*-Regular,BoldFont=*-Semibold,ItalicFont=*-It,BoldItalicFont=*-SemiboldIt]
\setmonofont{JetBrainsMono}[Path=fonts/,Extension=.ttf,UprightFont=*-Regular,Scale=MatchLowercase]
\setmathfont{Latin Modern Math}
\definecolor{accent}{HTML}{B3301B}
\definecolor{inksoft}{HTML}{625C51}
\hypersetup{colorlinks,urlcolor=accent,linkcolor=black}
\setlength{\parskip}{0.5em}
\setlength{\parindent}{0pt}

\newcounter{slide}
\newcommand{\slide}[3]{\stepcounter{slide}\subsection*{Slide \theslide: #1}{\color{inksoft}\small #2 \textperiodcentered\ #3}\par\vskip 0.3em}
\newcommand{\screen}[1]{\textbf{On screen.} #1\par}
\newcommand{\say}[1]{\textbf{Say.} #1\par}
\newcommand{\takeaway}[1]{\textbf{Takeaway.} #1\par}
\newcommand{\source}[1]{{\color{inksoft}\small Source: #1}\par}

\title{Speaker notes: transfinite game values in infinite chess}
\author{Sean Reid and Andrew Shargel}
\date{Draft of 21 October 2026}

\begin{document}
\maketitle

Each entry gives the presenter, a target time, what the slide shows, a script in full sentences, the one
thing the audience should keep, and where in the papers the claim lives. Total target: 25 minutes.
[EH14] is Evans and Hamkins, Integers 14 (2014) G2. [EHP17] is Evans, Hamkins and Perlmutter,
Integers 17 (2017) G4. Page numbers refer to the published PDFs.

\section*{Part A: ordinals and game values}

\slide{Title}{Sean}{0:20}
\screen{Title, our names, the journal.}
\say{We are presenting two papers from the journal Integers by Evans, Hamkins and Perlmutter about infinite
chess. The question they answer sounds like a joke and turns out to be real mathematics: if you are going to
lose a game of chess, how long can you make it last, and what kind of number measures that?}

\slide{Two papers, one topic}{Sean}{0:40}
\screen{The two citations and the split of the talk.}
\say{The 2014 paper builds the theory of game values and gives positions worth omega, omega squared and omega
cubed. The 2017 paper builds one position worth omega to the fourth, which is still the record for ordinary
infinite chess. I will do the theory and the small examples. Andrew will take you through the big position
and show it running.}

\slide{Infinite chess}{Sean}{0:50}
\screen{Five rules of the game.}
\say{Infinite chess is chess on the integer lattice. Same pieces, same moves, no edge anywhere. There is no
last rank, so pawns never promote. You win by checkmate, which happens at a definite move. A game that never
ends is a draw, and we drop the fifty-move and repetition rules, which only exist to force finite games.
Positions may have infinitely many pieces, as long as a rule says what is on each square. We assume you know
how the pieces move and nothing more.}
\source{[EH14] Section 1 and footnotes on pages 4 and 5; [EHP17] Section 2, page 4.}

\slide{Open games}{Sean}{1:30}
\screen{The definition of an open game.}
\say{Here is the one piece of general theory. A game is open for a player if whenever that player wins, the
win is already settled after finitely many moves. Chess is open for white, because checkmate is a fact about a
single position, and the same for black. A draw by infinite play is never settled at a finite stage, which is
why we will assign numbers from white's point of view and treat black separately. Everything I say next is
about open games in general. Chess is just where the examples come from.}
\takeaway{Open means a win shows up at a finite stage.}
\source{[EHP17] Section 2, first paragraph; [EH14] Section 2.}

\slide{Counting past the natural numbers}{Sean}{2:30}
\screen{The sequence of ordinals and the one fact.}
\say{We need numbers that go past the naturals. The ordinals do that. After every natural number comes omega.
After omega comes omega plus one, then omega plus two, and after all of those comes omega times two. Keep
going and you reach omega squared, omega cubed, omega to the fourth. Two kinds of ordinal matter: successors,
which have a predecessor, and limits like omega, which do not. You can add and multiply, and an expression like
omega cubed times five plus omega squared times seven plus omega plus fourteen is a single ordinal; you compare
two of them by the leading term first, like comparing numbers digit by digit. We need one fact: there is no
infinite strictly decreasing sequence of ordinals. The reason is that any nonempty collection of ordinals has a
least member, and a decreasing sequence would have no least member. Hold on to that fact; every proof today
rests on it.}
\takeaway{No infinite descent.}
\source{Standard; stated in [EH14] Section 2, page 5.}

\slide{A countdown game}{Sean}{2:00}
\screen{The table of starting ordinals and what black gets to do. The slide says this is our device.}
\say{Before any chess, here is a toy game we made up to show what the ordinals will mean. A counter shows an
ordinal. White wins when it reads zero. At a successor, white steps down by one. At a limit, black chooses any
smaller ordinal. Start at a natural number and black does nothing; the game lasts exactly that many moves.
Start at omega and black names a natural number, so the game lasts as long as black likes, but still finitely.
Start at omega times two and black gets to name a number twice. Start at omega squared and black names how many
times he will later get to name a number. White always wins by the no-descent fact. The ordinal measures how
much control black has over the length.}
\takeaway{The ordinal is a budget of announcements.}
\source{Our device. The announcement reading of values is [EH14] Section 2, page 5, and [EHP17] Section 2, page 5.}

\slide{Game value: the definition}{Sean}{2:00}
\screen{The three clauses of the definition.}
\say{Now the real definition, from white's point of view. A position where white has already won has value
zero. If it is white's turn, the value is alpha plus one where alpha is the least value white can move to. If
it is black's turn, the position has a value only if every legal black move leads to a valued position, and
then the value is the supremum of those. Read the quantifiers: white needs one good move, black needs all moves
to be covered. The valuation is the least assignment obeying these clauses, and some positions get nothing.
Finite value n is mate in n. Value omega means white wins and black's first move picks a natural number n with
the game lasting at least n more white moves. That is the countdown game inside chess.}
\takeaway{White: some move. Black: every move. Value is a supremum at black's turns.}
\source{[EH14] Definition 1, pages 4 and 5; [EHP17] Section 2, page 4.}

\slide{The fundamental observation}{Sean}{1:40}
\screen{Theorem and the first half of the proof.}
\say{Here is the theorem we prove in full. If a position has a value, white has a winning strategy. If it has
no value, black can prevent white from winning. First direction: white moves each turn to a position of least
value, which by the successor clause is one less. Black, by the supremum clause, can never move to a higher
value or to an unvalued position. So the values strictly decrease along the play, and by no-descent the play
reaches zero, which is a win.}
\takeaway{Reducing value is a winning strategy because the ordinals have no infinite descent.}
\source{[EH14] Section 2, page 5; [EHP17] Section 2, page 4.}

\slide{The fundamental observation, continued}{Sean}{1:20}
\screen{The second half of the proof and the Gale and Stewart line.}
\say{Second direction: suppose the position has no value. If it is white's turn, no white move reaches a
valued position, because one such move would have given the position a value by the successor clause. If it
is black's turn, some black move reaches an unvalued position, because otherwise the supremum clause would
have valued it. Black keeps choosing such moves and white never reaches zero. That is the whole proof. So the
valued positions are exactly the ones white can force to a win, and the same holds with the colours swapped.
Historically this argument is the proof of the Gale and Stewart theorem that open games are determined; we
only state that.}
\takeaway{Value exists if and only if white can force a win.}
\source{[EH14] Section 2, page 5; [EHP17] Section 2, page 4.}

\slide{Worked example: value omega}{Sean}{2:30}
\screen{Figure 3 of [EH14] and the main line.}
\say{A concrete position. Black to move, and the claim is value omega. The main line: black lifts the centre
rook as high as he likes, and white rolls the king upward with checks until it is mated against the rook. The
slide shows the line for height ten. The higher the rook, the longer the roll, so after black's first move
the values are finite but unbounded, and their supremum is omega. Two things must be checked. Any other black
first move loses at once; the slide lists the paper's refutations. And white has nothing faster: the tempting
skewer of the lifted rook takes the white rook off the f file, which lets the black king escape the net.
The paper gives these lines and says ``we claim''; we have not written out an exhaustive check, so this is a
proof sketch.}
\takeaway{One free choice by black, finite but unbounded, gives supremum omega.}
\source{[EH14] Figure 3 and the text on pages 8 and 9.}

\slide{Worked example: value omega squared}{Sean}{2:00}
\screen{Figure 7 of [EH14].}
\say{The next level nests the trick. An infinite pawn column splits the board. The marked pawn is a portcullis:
when it moves, the queens on the left flood the right half and mate. Black lifts the centre rook to height n, and
white must advance the column n times to free the bishop that keeps the door shut. But after every advance,
black's free rook on the right checks the white king from a distance k of black's choosing, and white has to walk
toward it for k moves before black must retreat and white can advance again. So black announces n, the number
of rounds, and in each round announces k. That is omega squared.}
\takeaway{Announce how many announcements: omega squared.}
\source{[EH14] Figure 7 and pages 13 to 15.}

\slide{From omega squared to omega cubed}{Sean}{1:30}
\screen{Figure 10 of [EH14].}
\say{One more nesting. The attacked bishop at m6 runs up the diagonal past as many towers as black likes
before it is captured. Each tower is a copy of the previous position: a rook that goes up, a pawn column that
marches, and the harassing rook on the left. The towers are consumed right to left. k towers give at least omega
squared times k, k is unbounded, so the value is omega cubed. Keep this picture; the omega to the fourth
position is this one with the harassing rook replaced by something stronger.}
\source{[EH14] Figure 10 and pages 17 to 19.}

\slide{What the 2014 paper proves and asks}{Sean}{1:30}
\screen{The list of results with their labels.}
\say{To close my half. The 2014 paper proves that positions with values omega, omega squared times k and omega
cubed exist. It defines the omega one of chess, the supremum of all values, and proves upper bounds we only
state: for finite positions the values are below the Church Kleene ordinal, and for arbitrary positions below
omega one. They conjecture these bounds are sharp. In three-dimensional infinite chess they prove every
countable ordinal occurs. Before 2015 the record in ordinary infinite chess was omega cubed. Andrew will show
you omega to the fourth.}
\source{[EH14] Theorem 4 page 10, Conjecture 5 page 26, Theorem 10 page 32; [EHP17] Section 1, page 2.}

\section*{Part B: the position with value omega to the fourth}

\slide{The theorem}{Andrew}{1:00}
\screen{Theorem 1 of [EHP17] and the meaning of its words.}
\say{Here is the result. There is a position in infinite chess with game value omega to the fourth; it is
computable and has infinitely many pieces. Position means an assignment of pieces to squares with black to
move. Computable means a program can tell you what is on any square. Value omega to the fourth means, in Sean's
countdown language, four nested kinds of announcement for black.}
\source{[EHP17] Theorem 1, page 3.}

\slide{The position}{Andrew}{1:30}
\screen{Figure 1 of [EHP17] and the three components.}
\say{This is the whole position. Three parts. The throne room at the upper left holds both kings, locked in.
The rook towers run off to the upper right; that is where white makes slow progress. The lower half is the
bishop cannon battery, which is black's way of forcing white to answer threats. Both the towers and the
battery continue forever.}
\source{[EHP17] Figure 1 and page 6.}

\slide{The throne room}{Andrew}{1:30}
\screen{Figure 2 of [EHP17] and the tempo lemma.}
\say{Nothing in the throne room can move, and both kings are exposed to a bishop of the right colour. A white
bishop mates at once from the marked diagonal. A black bishop that reaches its diagonal needs three more moves
to mate. That difference of three versus zero is the lemma everything else uses: if black's bishop reaches its
diagonal with a white guard two moves behind, white still mates first.}
\source{[EHP17] Figure 2 and pages 7 and 8.}

\slide{The rook towers}{Andrew}{2:00}
\screen{Figure 3 of [EHP17] and the four steps.}
\say{The towers are the omega cubed machine. Black's key bishop at f7 is attacked and runs up the diagonal as
far as black likes; that choice sets how many towers white must get through, the count r. White's pawn takes
it, pierces the wall, and a key pawn advances to attack the guard pawn of a tower. The attacked rook goes up as
high as black likes, the count h, and after it is captured white has to march the hole down one pawn move at a
time. When the hole arrives, bishops in the next tower are released and everything repeats one tower to the
left. The last tower frees the mating bishop.}
\source{[EHP17] Figure 3 and pages 8 to 10.}

\slide{A gateway terminal wing}{Andrew}{1:30}
\screen{Figure 4 of [EHP17].}
\say{A gateway wing is a long diagonal hallway with gates off it. A black bishop on the hallway can land on any
attack square, each of which threatens a door pawn. If white ignores the threat, black takes the door, exits,
and reaches the throne room three moves ahead of the guard, which mates by the lemma. So white must move the
guard one step, after which it follows two behind, which the lemma says is enough. A wing of length n gives
n plus one forced replies: the count g.}
\source{[EHP17] Figure 4 and pages 10 and 11.}

\slide{A bishop cannon}{Andrew}{1:30}
\screen{Figure 5 of [EHP17].}
\say{A cannon is a stack of bishops each capped by a pawn. Moving a capped bishop out uncovers a pawn threat
that white must meet, so each shot is a forced reply, and the fired bishop lands in a fresh wing and spends g
more forced replies there. Opening a cannon costs black one tempo, which is white's only free move, and buys
n shots: the count b. The cannons grow along the battery, so black can always open a bigger one.}
\source{[EHP17] Figure 5 and pages 11 to 13.}

\slide{The accounting}{Andrew}{2:00}
\screen{The ordinal with four terms and the three rules.}
\say{Here is the whole position as a single ordinal: omega cubed times the towers left, plus omega squared
times white's moves left in this tower, plus omega times the shots left, plus the forced replies left. A gateway
attack lowers g. A shot lowers b and raises g by a lot. Opening a cannon raises b by a lot but gives white a
free move, which lowers h. When h reaches zero the next tower starts, so r drops and black picks a new h. Black
can never raise a leading term, and white lowers one whenever black runs out of forced replies. The paper's
analysis of this is careful but written informally, so we label it a sketch.}
\source{[EHP17] Section 3, pages 13 to 15.}

\slide{Why deviations lose}{Andrew}{1:30}
\screen{Four one-line summaries.}
\say{Why can black not do better? Extra bishops are penned in and can only threaten door pawns, which white
answers as before. Wild rook moves in distant towers can be ignored a fixed fraction of the time, so they buy
less than the forced replies they replace. Exiting a gateway with the guard two behind loses the race. And white
cannot shortcut either: no bishop escapes the towers early, and capturing more pawns to the right only frees
bishops that are stuck above the pawn diagonal. Each line here stands for a paragraph of the paper.}
\source{[EHP17] pages 15 to 17.}

\slide{Live demo}{Andrew}{3:30}
\screen{The browser at 16:9, light theme, \url{https://infinite-chess.dwainosaur.com}. The slide says the app is
ours and not part of the papers.}
\say{This is our own implementation, written for this presentation. It plays white's strategy from the paper on
the compact version of the position from Section 4 and shows the value live.}
\takeaway{Each kind of black announcement moves one term of the ordinal, and white gets a free move only when
black opens a cannon.}

The readings below come from the engine on release 0.3.0. Walk the steps once before recording and update any
that differ. Read the value aloud each time it changes and say which term moved.

\textbf{Part one, from a fresh game.}
\begin{enumerate}
  \item Click New game. The panel reads $\omega^4$. Say: four nested announcements still to come.
  \item Click Hint. The key bishop lights up, its square is ringed, and buttons list landings: 1 for one tower,
    3 for two towers, 6 for three, and so on. Say: any distance is on the main line; every three squares is one
    more tower.
  \item Click the button \texttt{3 2 towers}. White captures at once. The panel reads
    $\omega^3\cdot 2+\omega^2\cdot 12$, with $r=2$ and $h=12$. Say: black announced $r$; $h$ is how many free
    moves white needs in this tower.
  \item Click Hint. The front bishop of cannon 1 lights up with landings 10 for wing 1, 18 for wing 2, and so on.
    Click \texttt{10 wing 1}. The panel reads $\omega^3\cdot 2+\omega^2\cdot 12+\omega+2$. Say: opening a cannon
    raised $b$ and $g$.
  \item White replies ``climbing the inner wall'' and the panel reads $\omega^3\cdot 2+\omega^2\cdot 11+\omega+2$.
    Say: that was white's one free move; $h$ went from 12 to 11. Opening a cannon is the only time white gets one.
  \item Click Hint, then the Move to button. The bishop attacks gate 1 of wing 1. White answers ``answering wing 1
    gate 1'' and the panel reads $\omega^3\cdot 2+\omega^2\cdot 11+\omega+1$. Say: $g$ dropped by one and nothing
    else moved; that is a forced reply.
  \item Click Hint, then the Move to button: the bishop takes the door pawn. The panel's $g$ shows ``open'' and
    the value reads $\omega^3\cdot 2+\omega^2\cdot 11+\omega\cdot 2$. Say: the bishop is in the doorway; white's
    guard settles in behind it.
  \item Click Hint. Shooter 1 of cannon 1 lights up with landings for fresh wings. Click \texttt{19 wing 2}.
    White replies ``shoring up cannon 1'' and the panel reads $\omega^3\cdot 2+\omega^2\cdot 11+7$, with $b=0$
    and $g=7$. Say: $b$ went to zero and $g$ jumped; the shot bought seven more forced replies.
  \item Click Hint and play two or three of the pawn steps. Each lowers $g$ by one while white keeps ``shoring
    up cannon 1''. Say: every black move here costs white a move; $h$ has not moved since step 5.
\end{enumerate}

\textbf{Part two, the rook step.} Open \href{https://infinite-chess.dwainosaur.com/\#g=zZJPSNNhGMef7\_f3bqPcCh1RnWaR\_YE2Ak-ydSiw2urWQZPqYhAUFQX25xBkMrc8LCFizToUJrRahQc1sOwQzCZpmXWwPKVQpEU0vHWIEbbe-P2kP2vt9Ht-vM\_z\_X6-z\_vubrDb9-0F7R63W0QoUl7hWU6IKCUCrJKqqvUu4eq1gm\_nc99cD7hmRXlFjS1Xlzldi\_JzO8qqqwMrQRFP5Ub\_j-d63wFbbe0xt3DLNnP9rZs8lfuXms3rOlFIXV07XML6BnOlXTs9lafWzaekK96Bq7Exhg2kiNd35NCvTOoKT-BoanqBJcLjJ82ZThz1-rpQ8zvSusV7IIIIPmOxMIo2mNucQwu8vj5s\_xOjn\_ZMlUAC7Vwg7MAVC8PLuAivL4M9f2OoG6dY1olOnOHm3IX4A924i0Ko6y5DtPWgB8\_pFvbinkW8PnTDH8jiYCEBdJB3lEEM4hNdwscYsgDJ4BH8gTBP\_wsQHajVsI9jHDHDIXyFCQug13gJfyDO8wW9m\_nJrhvOaUwjZUA4g48WZB\_wFv5AkteKQqYTpg1bFlnEjPrc2w2GWtjMYlLoNFOGhBnmjOEStjJK831F2MxgaJS3i7ovc-Kzyh5nnG3KIbzEDgviBC8wGJrkw\_9KrJNfVc4kk7yhILzJlAX5LXYxGJrlaEmQ6wkGVEU\_-5lWSnifAxYJHrCPwVDEmCypBHqSCbVsmMOcUguFI3xmkeQpMwyGEsZsSSbRE31Raoxj7FWHhW-YZu4HoozvtR2Yq0UR-Z58v2g1-BU}{the prepared link}, which loads the same game 606 plies in
(move 304), with cannon 12 spent and the first rook under attack. Say what it is: the same game, far along, saved
as a link.
\begin{enumerate}
  \setcounter{enumi}{9}
  \item The panel reads $\omega^3\cdot 2+\omega$, with $h=0$. Say: white has finished this tower's free moves;
    the next thing to happen is the second kind of announcement.
  \item Click Hint. The rook at $(47,57)$ lights up with an upward open line and the form prefilled. Type 5 and
    click Move. The panel reads $\omega^3\cdot 2+\omega^2\cdot 14+\omega$, with $h=14$. Say: black announced $h$;
    the higher the rook, the more pawn moves white needs.
  \item White captures the rook and the panel reads $\omega^3+\omega^2\cdot 13+\omega$, with $r=1$. Say: the
    tower is spent, $r$ dropped to 1, and the whole cycle starts again with one tower left.
\end{enumerate}

\textbf{If the app misbehaves.} If the panel does not update after white's move, wait two seconds; white runs in
a background worker and replays the game each move. If a reading differs from this list by one in $g$ or $b$,
carry on; say the term and the direction, not the number. If the link fails to load, play part one only and
describe the rook step from slide 17. Undo exists; use it if a click lands on the wrong square.
\source{Supplementary material; \url{https://github.com/sean-reid/infinite-chess}. The link above is the game's
move list in the URL; regenerate it from the engine if white's strategy changes.}

\slide{The compact battery}{Andrew}{0:40}
\screen{Figure 6 of [EHP17].}
\say{Section 4 folds the battery onto a quarter plane with a vertical wall. Same gateways, same accounting,
and a free quarter plane left over for bigger constructions. The demo uses this version.}
\source{[EHP17] Section 4, pages 21 and 22.}

\slide{Summary}{Andrew}{1:00}
\screen{Four lines with labels.}
\say{Open games have ordinal values, and a value exists exactly when white can force a win; that we proved.
Omega is one announcement, omega squared nests them, omega cubed nests again with towers; those we sketched from
the 2014 paper. Replacing the harassing rook by a cannon battery multiplies by omega squared instead of omega,
which gives omega to the fourth; that we sketched from the 2017 paper. Whether every countable ordinal is a
value in ordinary infinite chess is open.}

\slide{References and credits}{Andrew}{0:20}
\screen{Full citations, the DOIs, the figure credit, the repository.}
\say{The papers, the figure credit and our code are on this slide. Thank you.}

\section*{Glossary}

\begin{description}
  \item[Open game] A game in which any win for the designated player is already decided after finitely many moves.
  \item[Ordinal] A number in the well-ordered extension of the naturals: $0,1,2,\ldots,\omega,\ldots,\omega\cdot 2,\ldots,\omega^2,\ldots$
  \item[Successor, limit] $\alpha+1$ is the successor of $\alpha$; a limit ordinal such as $\omega$ has no predecessor.
  \item[Well-founded] No infinite strictly decreasing sequence exists. True of the ordinals.
  \item[Game value] The ordinal assigned to a position by the recursion on slide 7, from white's point of view.
  \item[Announcement] Black's move at a limit stage, choosing a smaller ordinal; in the countdown game, choosing a number.
  \item[Omega one of chess] The supremum of the game values of all positions; lightface over finite positions, boldface over all.
  \item[Church--Kleene ordinal] The least ordinal no computable well-ordering reaches; an upper bound for the values of finite positions.
  \item[Throne room] The locked chamber holding both kings in the $\omega^4$ position.
  \item[Rook tower] One repeating unit of white's main attack: guard pawn, rook, pawn column, bishops.
  \item[Gateway terminal wing] A diagonal hallway with gates, where a black bishop makes forced-reply threats.
  \item[Bishop cannon] A stack of pawn-capped bishops; firing one is a forced-reply threat.
  \item[$r$, $h$, $b$, $g$] Towers left, white moves left in the current tower, shots left in the active cannon, forced replies left in the active wing.
\end{description}

\end{document}
